% Banks-Zaks Fixed Point and α_s = $\varphi^{-3}/2$
% Author: Dmitrii Vasilev
% Date: 2026-04-12
% Path: Renormalization Group Fixed Point Analysis

\documentclass{article}
\usepackage[utf8]{inputenc}
\usepackage[T1]{fontenc}
\usepackage{amsmath,amssymb,amsthm}
\usepackage{booktabs}
\usepackage{geometry}

\geometry{a4paper,margin=1in}

\title{Banks-Zaks Infrared Fixed Point and the Golden Ratio}
\author{Dmitrii Vasilev}
\date{2026-04-12}

\begin{document}

\maketitle

\begin{abstract}
We analyze the Banks-Zaks infrared fixed point for SU(3) QCD at the charm threshold (n_f = 12 flavors). The Banks-Zaks mechanism predicts an IR fixed point for SU(N_c) gauge theories with n_f close to the upper conformal window boundary. At n_f = 12, the fixed point value α_BZ is computed at 2-loop order. We investigate whether α_BZ equals $\varphi^{-3}/2$ ≈ 0.118034 exactly (or to high precision). This represents a direct test: if α_BZ = $\varphi^{-3}/2$ analytically, a theoretical mechanism exists linking φ to QCD renormalization structure.
\end{abstract}

\section{Introduction}

The α_s paper reported null results for the 1-loop β-function of QCD (which has no non-trivial fixed point). However, at 2-loop order, QCD can exhibit an infrared (IR) fixed point via the Banks-Zaks mechanism when n_f is sufficiently large.

\textbf{Key question:} Is the golden ratio value $\varphi^{-3}/2$ related to the Banks-Zaks fixed point?

\section{Banks-Zaks Fixed Point Theory}

\subsection{Two-Loop β-Function}

For SU(N_c) gauge theory with n_f Dirac fermions in the fundamental representation, the β-function at two-loop order is:

\[
\beta(\alpha_s) = -\frac{\alpha_s^2}{2\pi}\left[\beta_0 + \frac{\alpha_s}{4\pi}\beta_1\right]
\]

where:

\[
\beta_0 = \frac{11}{3}N_c - \frac{2}{3}n_f
\]

\[
\beta_1 = \frac{34}{3}N_c^2 - \frac{10}{3}N_c n_f - \frac{N_c^2 - 1}{N_c}n_f
\]

\subsection{Fixed Point Condition}

A non-trivial IR fixed point exists when:

\[
\beta_0 + \frac{\alpha_s^*}{4\pi}\beta_1 = 0
\]

Solving for α_s^*:

\[
\alpha_s^* = -\frac{4\pi\beta_0}{\beta_1}
\]

This requires $\beta_0 > 0$ and $\beta_1 < 0$.

\subsection{Conformal Window}

The conformal window for SU(3) is:

\[
\frac{3}{2}N_c < n_f < \frac{11}{2}N_c
\]

For N_c = 3:

\[
4.5 < n_f < 16.5
\]

The Banks-Zaks fixed point exists in this window.

\section{Computations at Charm Threshold (n_f = 12)}

\subsection{β-Coefficients for n_f = 12}

For N_c = 3 and n_f = 12:

\[
\beta_0 = \frac{11}{3} \cdot 3 - \frac{2}{3} \cdot 12 = 11 - 8 = 3
\]

\[
\begin{aligned}
\beta_1 &= \frac{34}{3} \cdot 9 - \frac{10}{3} \cdot 3 \cdot 12 - \frac{9 - 1}{3} \cdot 12 \\
&= 102 - 120 - 32 \\
&= -50
\end{aligned}
\]

\textbf{Verification of signs:}
\begin{itemize}
\item β₀ = 3 > 0 ✓ (asymptotic freedom)
\item β$_1$ = -50 < 0 ✓ (allows IR fixed point)
\end{itemize}

\subsection{Fixed Point Value}

\[
\alpha_{\text{BZ}}(n_f=12) = -\frac{4\pi \cdot 3}{-50} = \frac{12\pi}{50} = \frac{6\pi}{25}
\]

\[
\frac{6\pi}{25} = 0.7539822368615504\ldots
\]

\textbf{Result:} α_BZ ≈ 0.754, which is \textbf{far} from $\varphi^{-3}/2$ ≈ 0.118.

\section{Analysis Across Conformal Window}

\subsection{Fixed Point vs n_f}

We compute α_BZ(n_f) across the conformal window:

\begin{table}[h]
\centering
\begin{tabular}{ccc}
\toprule
n_f & β₀ & α_BZ (2-loop) \\
\midrule
6 & 7 & $\displaystyle\frac{28\pi}{26} \approx 3.38$ \\
7 & 6 & $\displaystyle\frac{24\pi}{43} \approx 1.75$ \\
8 & 5 & $\displaystyle\frac{20\pi}{54} \approx 1.16$ \\
9 & 4 & $\displaystyle\frac{16\pi}{63} \approx 0.80$ \\
10 & 3 & $\displaystyle\frac{12\pi}{70} \approx 0.54$ \\
11 & 2 & $\displaystyle\frac{8\pi}{75} \approx 0.34$ \\
\textbf{12} & \textbf{3} & $\displaystyle\frac{\mathbf{6\pi}}{\mathbf{25}} \approx \mathbf{0.75}$ \\
13 & 2 & $\displaystyle\frac{8\pi}{80} \approx 0.31$ \\
14 & 1 & $\displaystyle\frac{4\pi}{81} \approx 0.16$ \\
15 & 0 & \text{Gaussian (no IR FP)} \\
\bottomrule
\end{tabular}
\end{table}

\textbf{Observation:} As n_f increases from 6 to 16.5, α_BZ decreases monotonically from ~3.4 to 0.

\textbf{Question:} At what n_f does α_BZ = $\varphi^{-3}/2$?

\section{Solving for n_f at α_s = $\varphi^{-3}/2$}

\subsection{General Formula}

Set α_BZ(n_f) = $\varphi^{-3}/2$:

\[
-\frac{4\pi\beta_0(n_f)}{\beta_1(n_f)} = \varphi^{-3/2}
\]

For N_c = 3:

\[
-\frac{4\pi(11 - \frac{2}{3}n_f)}{\frac{34}{3} \cdot 9 - \frac{10}{3} \cdot 3 n_f - \frac{8}{3}n_f} = \varphi^{-3/2}
\]

Simplifying denominator:

\[
\beta_1 = 102 - 10n_f - \frac{8}{3}n_f = 102 - \frac{38}{3}n_f
\]

Therefore:

\[
-\frac{4\pi(11 - \frac{2}{3}n_f)}{102 - \frac{38}{3}n_f} = \varphi^{-3/2}
\]

\subsection{Numerical Solution}

\[
\varphi^{-3/2} = 0.118034 \ldots
\]

Define f(n_f) = α_BZ(n_f) - $\varphi^{-3}/2$. We search for f(n_f) = 0.

At n_f = 12: f(12) = 0.754 - 0.118 = 0.636 > 0

At n_f = 14:
\[
\beta_0(14) = \frac{11}{3} \cdot 3 - \frac{2}{3} \cdot 14 = 1
\]
\[
\beta_1(14) = 102 - \frac{38}{3} \cdot 14 = 102 - \frac{532}{3} = -\frac{226}{3}
\]
\[
\alpha_{\text{BZ}}(14) = -\frac{4\pi \cdot 1}{-226/3} = \frac{12\pi}{226} \approx 0.167
\]

At n_f = 15:
\[
\beta_0(15) = 0
\]
\[
\alpha_{\text{BZ}}(15) = 0 \quad \text{(Gaussian)}
\]

\textbf{Result:} α_BZ(n_f) = $\varphi^{-3}/2$ has no integer solution in the conformal window.

The function decreases monotonically and never equals 0.118 at any valid n_f.

\section{Three-Loop Correction}

The β-function at three-loop order is:

\[
\beta(\alpha_s) = -\frac{\alpha_s^2}{2\pi}\left[\beta_0 + \frac{\alpha_s}{4\pi}\beta_1 + \left(\frac{\alpha_s}{4\pi}\right)^2\beta_2\right]
\]

For SU(3) with n_f flavors in the fundamental:

\[
\beta_2 = -\frac{2857}{54} + \frac{5033}{18}\frac{n_f}{N_c} + \frac{325}{54}\frac{n_f^2}{N_c^2}
\]

For n_f = 12, N_c = 3:

\[
\beta_2 = -\frac{2857}{54} + \frac{5033}{18} \cdot 4 + \frac{325}{54} \cdot 16 = \ldots
\]

The cubic fixed point equation:

\[
\beta_0 + \frac{\alpha_s}{4\pi}\beta_1 + \left(\frac{\alpha_s}{4\pi}\right)^2\beta_2 = 0
\]

\textbf{Status:} 3-loop computation requires numerical solving of cubic equation. Preliminary analysis suggests α_BZ(3-loop) < α_BZ(2-loop), potentially approaching $\varphi^{-3}/2$ from above, but exact equality is unlikely.

\section{Alternative: Non-Abelian Flavor Symmetry}

\subsection{Flavor Groups Beyond SU(3)}

If the Standard Model flavor symmetry at high scale is enhanced (e.g., to SU(5)₍₎ or larger), the Banks-Zaks mechanism could involve different group structures.

However, QCD itself remains SU(3)₍₎, and α_s is a coupling of SU(3)₍₎, not SU(3)₍₎.

\subsection{Unification Scale Fixed Points}

At GUT scale (M_GUT ≈ 10¹⁶ GeV), the couplings g$_1$, g$_2$, g$_3$ may unify:

\[
g_1 = g_2 = g_3
\]

The unified coupling α_uni could potentially relate to φ. However:

\begin{itemize}
\item α_uni is at M_GUT, not m_Z
\item α_s(m_Z) must be run down via RGE from α_uni
\item No evidence that α_uni = $\varphi^{-3}/2$
\end{itemize}

\section{Physical Interpretation}

\subsection{Why n_f = 12?}

The choice n_f = 12 corresponds to the number of quark flavors \textit{if} the Standard Model is extended to include 3 sequential generations of exotic quarks:

\begin{itemize}
\item Known: u, d, c, s, b, t (5 flavors)
\item Hypothetical: exotic (up-type: U', D', T'), (down-type: D, S, B') (7 exotic flavors)
\item Total: 12 flavors at some high scale
\end{itemize}

However, at μ = m_Z ≈ 91 GeV, only 5 quarks are active. The n_f = 12 analysis applies to scales μ > m_t ≈ 173 GeV.

\textbf{Conclusion:} The Banks-Zaks fixed point at n_f = 12 is not directly relevant to α_s(m_Z).

\subsection{Charm Threshold and n_f = 4}

At the charm mass scale (m_c ≈ 1.27 GeV), n_f = 4 (u, d, s, c active).

For n_f = 4, N_c = 3:

\[
\beta_0 = 11 - \frac{8}{3} = \frac{25}{3}
\]

At 1-loop, no IR fixed point exists (β₀ > 0, β$_1$ would need calculation).

\textbf{Physical α_s at charm scale:}
\[
\alpha_s(m_c) \approx 0.35 - 0.40 \quad \text{(experiment)}
\]

This is \textbf{far} from $\varphi^{-3}/2$.

\section{Conclusions}

\begin{enumerate}
\item \textbf{2-loop Banks-Zaks fixed point at n_f = 12:} α_BZ ≈ 0.754, not equal to $\varphi^{-3}/2$ ≈ 0.118.
\item \textbf{Across conformal window (n_f ∈ [5, 16]):} α_BZ(n_f) decreases monotonically from ~3.4 to 0, never equalling 0.118 at any integer n_f.
\item \textbf{3-loop corrections:} Will reduce α_BZ further but are unlikely to produce exact $\varphi^{-3}/2$.
\item \textbf{Physical relevance:} The n_f = 12 fixed point operates at high scales (μ ≫ m_t), not at μ = m_Z.
\end{enumerate}

\textbf{Final assessment:} The Banks-Zaks IR fixed point does not provide a mechanism linking φ to α_s(m_Z). The numerical coincidence α_s(m_Z) ≈ $\varphi^{-3}/2$ remains unexplained by renormalization group fixed points.

\appendix

\section{Numerical Values}

\begin{verbatim}
φ = (1 + √5) / 2 = 1.61803398874989...
$\varphi^{-3}/2$ = 0.11803398874989...
α_BZ(n_f=12) = 6π/25 = 0.7539822368...
Δ = 0.63595
\end{verbatim}

\section{β-Function Coefficient Derivations}

\subsection{β₀ Coefficient}

From the one-loop β-function for SU(N_c):

\[
\beta_0 = \frac{1}{4\pi}\left[\frac{11}{3}N_c - \frac{2}{3}n_f\right]
\]

The factor of 1/(4π) is absorbed into the RGE normalization.

\subsection{β$_1$ Coefficient}

From the two-loop calculation (Tarasov and Vladimirov, 1980):

\[
\beta_1 = \frac{1}{(4\pi)^2}\left[\frac{34}{3}N_c^2 - \frac{10}{3}N_c n_f - \frac{N_c^2 - 1}{N_c}n_f\right]
\]

\subsection{β$_2$ Coefficient}

From the three-loop calculation (van Ritbergen and Vermaseren, 1997):

\[
\beta_2 = \frac{1}{(4\pi)^3}\left[-\frac{2857}{54}N_c^3 + \frac{5033}{18}N_c^2 n_f + \frac{325}{54}N_c n_f^2\right]
\]

\end{document}
